\FloatBarrier
\section{Related work and what is new}\label{sec:related}

\paragraph{Forecast evaluation.} Comparing the out-of-sample accuracy of competing forecasters on the same realisations, rather than their in-sample fit, is the established discipline for the question asked here \citep{Diebold1995,HyndmanKoehler2006,HyndmanAthanasopoulos2021}, and rolling-origin evaluation is its standard form for time series \citep{Tashman2000,Bergmeir2012}. The hazard that discipline guards against is selection after the fact: when many candidates are scored on one realisation, the best observed score is an upward-biased estimate of skill and per-comparison tests do not restore validity. White's reality check constructs the distribution of the best statistic across candidates under the null that none has skill \citep{White2000}, and stepwise procedures formalise the correction with control of the familywise error rate \citep{RomanoWolf2005}. The present design addresses the same hazard by a more restrictive and less general route, and the difference should be stated. Those procedures take a set of already-computed scores and ask which survive correction for having been selected. Here the model classes, baselines and scoring cores were declared for each system; the Edwards H1/H2 gates were frozen before their respective scores, while the cod 5\% band and some comparator declarations were formalised after the first pass. These completions are disclosed, and the operating characteristics are estimated by simulation conditional on the cod rule as applied. This design is not a familywise-error correction and is not presented as one.

\paragraph{Simple benchmarks are hard to beat.} That naive and simple statistical methods are close to the most accurate, and that sophisticated methods frequently fail to beat them, is the most consistently replicated result in forecasting: it organises the M3 and M4 competitions \citep{Makridakis2000,Makridakis2018}, in the latter of which none of the pure machine-learning entries was more accurate than the combination benchmark, and it has been found on fish population series \citep{Ward2014} and on simulated ecological dynamics \citep{Perretti2013}. The two results here are single-system instances of that phenomenon, and the scope difference is material: the competitions aggregate over tens of thousands of series and dozens of methods, whereas this paper scores two ladders on two systems. A cross-series result about method classes entails nothing about a particular stock or basin, and a two-system result does not generalise to method classes. What the two kinds of result share is the test, and what this paper adds to the phenomenon is a reason, in equation~\eqref{eq:decomp}, why a mechanistic increment forecaster loses to persistence when it does.

\paragraph{Fisheries.} Structural claims in stock assessment are usually examined through fit, retrospective pattern and simulation self-consistency \citep{HilbornWalters1992,Punt2016,Maunder2015}; where prediction skill is examined directly, it is computed on the fitted survey or catch-rate index through hindcast cross-validation \citep{Kell2016,Kell2021,Carvalho2021,Kokkalis2024}, while advice concerns the state, and the difficulties of fitting surplus-production models to series with little contrast are long documented \citep{HilbornWalters1992,LudwigWalters1985,Polacheck1993,Punt2003,PedersenBerg2017}. The Northern cod collapse and its non-recovery have an extensive literature establishing the empirical facts \citep{Hutchings1994,Myers1997,Walters1996,Hutchings2005,RoseRowe2015,RoseWalters2019,Rose2026} and debating depensation \citep{Shelton1999,Hutchings1996,Myers1995,Liermann1997}. This paper's contribution to that literature is different in kind: it measures what those facts do to the forecast skill of a familiar model class from fixed origins, locates the deficit in the collapse transitions, proves that those transitions lie outside the class at the recorded catch, and reports the power of the test. It is a conditional hindcast, scored against a retrospective vintage with catch supplied along the horizon; it is not an operational forecast evaluation and licenses no claim about the performance of the assessment process as it ran.

\paragraph{Groundwater.} Groundwater-level forecasting at index wells is a mature area with statistical, neural-network and process-based approaches on annual and sub-annual series \citep{Daliakopoulos2005,AdamowskiChan2011}, and it has begun to institutionalise benchmark culture: GEMS-GER standardises 32 years of weekly observations from 3\,207 German wells with three benchmark models of increasing complexity \citep{Ohmer2026}, and a systematic comparison of nine machine-learning architectures on a karst catchment finds that aquifer complexity rather than algorithm choice bounds predictability \citep{Zhu2026}. Those studies compare model families against one another. None subjects a deliberately simple process-based water balance to a scored ablation against persistence and the training mean under a dated pre-score computational gate, and none reports a training mean beating persistence at a multi-year horizon on this basin. On the Edwards Aquifer, regional flow has long been represented, and debated, through equivalent-porous-media and lumped approaches \citep{Scanlon2003}; the one-pool map is not a karst model and the lumped-versus-distributed question is inherited rather than settled here. The climate predictors enter because ENSO phases shift North American precipitation \citep{Ropelewski1986}, and the finding that they add at most 0.13\,ft at the annual origin is a statement about this information set, not about ENSO.

\paragraph{What is new.} Relative to the two single-system studies and their earlier two-part combination, this article contributes (a) one scored test reported on two systems under disclosed, system-specific point rules without pooling, with the Edwards score gates and class dispositions separated; (b) Proposition~\ref{prop:increment}, an exact one-step MSE identity separating driver substitution from residual covariance, with a half-autocorrelation threshold only under the stated stationary, mean-zero and orthogonality conditions, and a constant-increment corollary; (c) Remark~\ref{rem:crossover}, the half-autocorrelation crossover of persistence and the training mean, evaluated on J-17 origins and shown not to describe the non-stationary cod series; (d) the correction of the earlier cross-system reading, which attributed the cod failure to a weakly autocorrelated driver; and (e) for the cod half, the period decomposition, identification result and operating characteristics, reported alongside—but not equated with—the J-17 realized-flux conditional diagnostic, which uses unavailable future drivers and is not an operating-characteristic estimate.
